\documentclass[10pt]{article}
\usepackage[T1]{fontenc}
\usepackage[utf8]{inputenc}
\usepackage[french]{babel}
\usepackage{amsmath,amssymb,mathrsfs,mathpazo}
\usepackage{geometry}
\geometry{a4paper,margin=2cm}
\usepackage{graphicx,array,tabularx,tabularray,colortbl,multirow,multicol,booktabs,enumitem,pifont,ulem,fancybox,diagbox,cancel}
\usepackage{pstricks,pstricks-add,pst-plot,pst-tree,pst-node,pst-text,pst-func,pst-math,pst-bspline,pst-3dplot,pst-eucl,pst-solides3d}
\let\pscalculate\undefined
\usepackage{pst-bezier}
\usepackage{tikz,pgfplots,tkz-tab}
\usetikzlibrary{arrows,arrows.meta,patterns,calc,positioning,trees,decorations.markings,shapes.geometric,shapes.misc,shapes.arrows,intersections,decorations.pathmorphing,decorations.pathreplacing,backgrounds,shadows,babel}
\pgfplotsset{compat=1.18}
\usepgfplotslibrary{fillbetween}
\usepackage[np]{numprint}
\usepackage{listings,makecell,pas-tableur,twoopt}
\let\barre\undefined
\usepackage[pstricks]{tablvar}
\definecolor{darkgreen}{RGB}{0,100,0}
\providecommand{\red}{\color{red}}
\providecommand{\blue}{\color{blue}}
\providecommand{\green}{\color{green}}
\providecommand{\black}{\color{black}}
\providecommand{\white}{\color{white}}
\providecommand{\gray}{\color{gray}}
\providecommand{\cyan}{\color{cyan}}
\providecommand{\magenta}{\color{magenta}}
\providecommand{\yellow}{\color{yellow}}
\providecommand{\violet}{\color{violet}}
\providecommand{\orange}{\color{orange}}
\providecommand{\euro}{\texteuro}
\providecommand{\EUR}{\texteuro}
\providecommand{\ieme}{\textsuperscript{e}}
\providecommand{\degres}{\ensuremath{^\circ}}
\providecommand{\degree}{\ensuremath{^\circ}}
\providecommand{\vv}[1]{\overrightarrow{#1}}
\providecommand{\Coord}[1]{#1}
\providecommand{\decosix}{}
\providecommand{\starredbullet}{\ensuremath{\bullet}}
\pagestyle{empty}
\setlength{\parindent}{0pt}
\usepackage{adjustbox}\DeclareUnicodeCharacter{2236}{:}\DeclareUnicodeCharacter{2190}{\ensuremath{\leftarrow}}
\providecommand{\R}{\mathbb{R}}
\providecommand{\N}{\mathbb{N}}
\providecommand{\Z}{\mathbb{Z}}
\providecommand{\Q}{\mathbb{Q}}
\providecommand{\D}{\mathbb{D}}
\providecommand{\C}{\mathbb{C}}
\providecommand{\e}{{\mathrm{e}}}
\providecommand{\vect}[1]{\overrightarrow{#1}}
\providecommand{\vectt}[1]{\overrightarrow{\mathrm{#1}}}
\providecommand{\vvt}[1]{\overrightarrow{\mathrm{#1}}}
\providecommand{\np}[2][]{{#2~#1}}
\providecommand{\vv}[1]{{\overrightarrow{#1}}}
\renewcommand{\d}{{\mathrm{d}}}
\providecommand{\degres}{{^\circ}}
\providecommand{\degre}{{^\circ}}
\providecommand{\blue}{\color{blue}}
\providecommand{\red}{\color{red}}
\providecommand{\enskip}{\quad}
\providecommand{\llbracket}{\lbrack\!\lbrack}
\providecommand{\rrbracket}{\rbrack\!\rbrack}
\providecommand{\euro}{{\text{€}}}
\providecommand{\no}{{\text{n°}}}
\providecommand{\interval}[2]{{[#1;#2]}}
\providecommand{\intervFO}[2]{{[#1;#2[}}
\providecommand{\dotfill}{\ldots}

\providecommand{\ds}{\displaystyle}
\providecommand{\textding}[1]{#1}
\providecommand{\Oij}{\(\mathrm{O};\vec{\imath},\vec{\jmath}\)}
\providecommand{\Oijk}{\(\mathrm{O};\vec{\imath},\vec{\jmath},\vec{k}\)}
\providecommand{\Ouv}{\(\mathrm{O};\vec{u},\vec{v}\)}
\providecommand{\intervFF}[2]{\ensuremath{[#1;#2]}}
\renewcommand{\euro}{\ifmmode\text{\texteuro}\else\texteuro\fi}
\DeclareUnicodeCharacter{00B0}{\ifmmode^{\circ}\else\textdegree\fi}
\DeclareUnicodeCharacter{20AC}{\ifmmode\text{\texteuro}\else\texteuro\fi}
% Énoncés : Éducation nationale. Transcription et corrigés : APMEP, Denis Vergès et contributeurs. Découpage et classement : CoopMaths, Sébastien Lozano, Éric Elter et contributeurs. Adaptation de la présentation : Maths Explorer.
% CC BY-SA : https://coopmaths.fr/www/about/
% https://coopmaths.fr/alea/static/eam/2026/tex/eam_specialite_2026_06_metropole_automatismes_cor.tex
\begin{document}
\noindent\textit{APMEP / CoopMaths. CC BY-SA. Présentation adaptée par Maths Explorer.}\par\medskip
\textbf{Pour cette première partie, aucune justification n'est demandée et une seule réponse est possible par question. Pour chaque question, reportez son numéro sur votre copie et indiquez votre réponse.}

	Une réponse fausse ou l'absence de réponse n'enlève aucun point.


	\vspace{0,25cm}
\subsubsection*{Question 1}
La forme développée de l'expression \mbox{$(3x-2)^2$} est :
		\begin{enumerate}[label=\alph*)]
\begin{multicols}{2}
			\item \mbox{$9x^2-4$}
			\item \mbox{$3x^2-12x+4$}

\end{multicols}
\begin{multicols}{2}
	\item \mbox{$9x^2-12x+4$}
	\item \mbox{$6x-4$}

	\end{multicols}
		\end{enumerate}


 \textcolor{red}{\textsc{Correction}}


	C'est une identité remarquable : \quad \mbox{$(3x-2)^2 = (3x)^2 - 2\times (3x) \times 2 + 2^2 = 9x^2-12x+4$}

	\hfill \framebox[1.2\width]{réponse C}

		\subsubsection*{Question 2}
\begin{minipage}[t]{0.45\linewidth}\vspace{0mm}
	Dans le repère (O, I, J) ci-contre, la droite \mbox{$(\Delta)$} a pour équation :
		\begin{enumerate}[label=\alph*)]
			\begin{multicols}{2}
				\item \mbox{$y=2x+2$}
				\item \mbox{$y=-2x+2$}

				\end{multicols}
			\begin{multicols}{2}
				\item \mbox{$y=-x+2$}
				\item \mbox{$y=x+2$}

				\end{multicols}
		\end{enumerate}

\end{minipage}\hspace{0.05\linewidth}\begin{minipage}[t]{0.45\linewidth}\vspace{0mm}
\begin{tikzpicture}[xscale=1,yscale=1,rotate=0,every node/.style={scale=1,rotate=0},baseline={(0,2.8)}]
	\clip(-1.6,-1.4)rectangle(3.6,3.2);
	\draw[blue!20, very thin](-1.6,-1.4)grid[xstep=0.2,ystep=0.2](3.6,3.2);
	\draw[blue!30,thick](-1.6,-1.4)grid(3.6,3.2);
	\draw[line width=0.3mm,->, >=triangle 45](-1.6,0)--(3.6,0);
	\draw[line width=0.3mm,->, >=triangle 45](0,-1.4)--(0,3.2);
	\draw[thick,shift={(0,0)}](2pt,0)--(-2pt,0) (0,2pt)--(0,-2pt) node[below left]{O};
	\draw[thick,shift={(1,0)}](2pt,0)--(-2pt,0) (0,2pt)--(0,-2pt) node[below]{I};
	\draw[thick,shift={(0,1)}](2pt,0)--(-2pt,0) (0,2pt)--(0,-2pt) node[left]{J};
	\draw[red, line width=0.3mm](-1.6,3.6)--(3.6,-1.6)node[pos=0.2,below left]{$(\Delta)$};
\end{tikzpicture}

\end{minipage}

 \textcolor{red}{\textsc{Correction}}

 L'ordonnée à l'origine est $2$.

 $(\Delta)$ passe par les points A\mbox{$(0;2)$} et B\mbox{$(2;0)$}, donc le coefficient directeur est donné par :

 $\dfrac{y_\mathrm{B} - y_{\mathrm{A}}}{x_{\mathrm{B}}-x_{\mathrm{A}}} = \dfrac{0 - 2}{2 - 0} = \dfrac{-2}{2} = -1$.

 L'équation réduite de $(\Delta)$ est \mbox{$y=-x+2$}

 \hfill \framebox[1.2\width]{réponse C}

\subsubsection*{Question 3}
Dans une classe de première, 75\% des élèves étudient le grec.

Les autres élèves étudient le latin : ils sont 9.

Le nombre d'élèves de cette classe de première est égal à :
\begin{enumerate}[label=\alph*)]
	\begin{multicols}{4}
		\item \mbox{$24$}
		\item \mbox{$30$}
		\item \mbox{$34$}
		\item \mbox{$36$}

		\end{multicols}
\end{enumerate}


 \textcolor{red}{\textsc{Correction}}


 Donc 25\% soit $\dfrac14$ des élèves étudient le latin, l'effectif de la classe est donc \quad \mbox{$9\times4=36$},

 \hfill \framebox[1.2\width]{réponse D}
\subsubsection*{Question 4}
Le prix d'un article augmente de 15\%.

Cela signifie que le prix de l'article a été multiplié par :
\begin{enumerate}[label=\alph*)]
	\begin{multicols}{4}
		\item \mbox{$\dfrac{15}{100}$}
		\item \mbox{$\np{1.15}$}
		\item \mbox{$\np{0.85}$}
		\item \mbox{$\np{1.115}$}

		\end{multicols}
\end{enumerate}


 \textcolor{red}{\textsc{Correction}}


 Augmenter une valeur de 15\% revient à la multiplier par \quad \mbox{$\dfrac{100+15}{100}=\np{1.15}$},

 \hfill \framebox[1.2\width]{ réponse B}
\subsubsection*{Question 5}
Parmi les réponses proposées, la valeur la plus proche de \mbox{$\dfrac{\np{150000}}{\np{3200}}$} est :
\begin{enumerate}[label=\alph*)]
	\begin{multicols}{4}
		\item \mbox{$5$}
		\item \mbox{$50$}
		\item \mbox{$500$}
		\item \mbox{$\np{5000}$}

		\end{multicols}
\end{enumerate}

 \textcolor{red}{\textsc{Correction}}


 \mbox{$\np{3200}\approx\np{3000}$}, donc \mbox{$\dfrac{\np{150000}}{\np{3200}}\approx\dfrac{150}{3}=50$},

 \hfill \framebox[1.2\width]{réponse B}
\subsubsection*{Question 6}
Une vidéo, d'une durée de 1 minute et 40 secondes, contient \np{2400}~images.

Le nombre d'images par seconde est égal à :
		\begin{enumerate}[label=\alph*)]
	\begin{multicols}{2}
		\item\np{60} images/seconde
		\item\np{24} images/seconde

		\end{multicols}
	\begin{multicols}{2}
		\item\np{120} images/seconde
		\item\np{15} images/seconde

		\end{multicols}
\end{enumerate}

 \textcolor{red}{\textsc{Correction}}


 1 minute et 40 secondes c'est 100 secondes, il y a donc 24 images par seconde,

 \hfill \framebox[1.2\width]{réponse B}
\subsubsection*{Question 7}
On considère une fonction $f$ définie sur $\mathbb{R}$ par :

\begin{center}
	\mbox{$f(x)=\np{0.5}(x-3)^2+10$}
\end{center}


On note \mbox{$\mathcal{C}$} sa courbe représentative dans un repère.

Un seul des quatre points ci-dessous appartient à la courbe \mbox{$\mathcal{C}$}. Lequel ?
\begin{enumerate}[label=\alph*)]
	\begin{multicols}{4}
		\item A~\mbox{$(-3;10)$}
		\item B~\mbox{$(3;\np{10.5})$}
		\item C~\mbox{$(3;10)$}
		\item D~\mbox{$(0;\np{19.5})$}
		\end{multicols}
\end{enumerate}

 \textcolor{red}{\textsc{Correction}}


 On a la forme canonique d'une fonction du second degré, dont le point C est clairement le sommet,

 \hfill \framebox[1.2\width]{réponse C}

\subsubsection*{Question 8}
On considère le nombre \mbox{$A=\dfrac{10^{201}\times10^{-4}}{(10^2)^{100}}$}

On peut affirmer que :
\begin{enumerate}[label=\alph*)]
	\begin{multicols}{4}
		\item \mbox{$A=-\np{0.001}$}
		\item \mbox{$A=\np{0.0001}$}
		\item \mbox{$A=\np{0.001}$}
		\item \mbox{$A=\np{1000}$}
		\end{multicols}
\end{enumerate}


 \textcolor{red}{\textsc{Correction}}

 \mbox{$\left(10^2\right)^{100}=10^{200}$}, \quad donc\quad  \mbox{$A=10^{201-4-200}=10^{-3}=\np{0.001}$},

 \hfill\framebox[1.2\width]{réponse C}

\end{document}
